% =====================================================================
\chapter{Kết luận}
\label{chap:ketluan}
% =====================================================================

\section{Kết quả đạt được}

Đề tài đã nghiên cứu và tái hiện đầy đủ bài báo ``\BaiBaoGoc'' \cite{raut2026hybrid} từ MATLAB sang
Python. Các kết quả chính gồm:
\begin{itemize}
  \item \textbf{Về lý thuyết}: hệ thống hóa các kiến thức nền tảng của đề tài, gồm mã khối 3DES và
        AES ở chế độ CBC, đệm PKCS\#7, hệ hỗn loạn Logistic, mã dòng XOR, giấu tin LSB đa bit, và các
        tiêu chí đánh giá (MSE, PSNR, SSIM, NPCR, UACI, entropy, hệ số tương quan). Từ đó em rút ra
        công thức sai số kỳ vọng $(L^2-1)/6$ của phép nhúng $k$-LSB, giải thích được vì sao bài báo
        đạt $\mathrm{MSE} \approx 0.5$ và PSNR $\approx \SI{51.15}{dB}$.
  \item \textbf{Về cài đặt}: xây dựng thư viện lõi gồm ba thuật toán mã hóa theo mẫu Strategy, lớp
        giấu tin LSB 1--3 bit vectơ hóa hoàn toàn bằng NumPy, header 16 byte tự mô tả cho phép bên
        nhận giải mã chỉ từ ảnh stego và khóa, và bộ độ đo đánh giá kép. Hệ thống có 347 ca kiểm thử
        tự động, tất cả đều đạt.
  \item \textbf{Về tái hiện}: xác định và giải quyết nhất quán mười điểm mù kỹ thuật của bài báo.
        Trên toàn bộ 8000 ảnh Linnaeus 5 (\num{24000} lượt chạy), \textbf{100\%} ảnh bí mật được khôi
        phục chính xác từng bit. MSE ($0.4956$), PSNR (\SI{51.18}{dB}) và NPCR ($\approx 99.6\%$) khớp
        sát Bảng~2; thứ tự tốc độ Chaotic $<$ AES $<$ 3DES được giữ nguyên. Thời gian 3DES trên ảnh
        $246\times300$ (\SI{3588}{\micro s}) gần như trùng số liệu của bài báo (\SI{3552}{\micro s}).
  \item \textbf{Về phân tích}: giải thích được các con số ``bất thường'' trong Bảng~2. UACI thấp hơn
        $33.46\%$ là do phân bố mức xám của ảnh bí mật (giá trị kỳ vọng tính được là $30.41\%$);
        entropy $\approx 7.05$ thực chất là entropy của ảnh stego, trong khi entropy bản mã đạt
        $7.993$.
  \item \textbf{Về đánh giá an toàn mở rộng}: chỉ ra bằng thực nghiệm các hạn chế mà bài báo chưa
        nêu. Lớp Chaotic + PRNG bị phá hoàn toàn bởi tấn công biết bản rõ và không có tính khuếch tán
        (đổi một điểm ảnh chỉ đổi một byte bản mã). CBC với IV cố định chỉ khuếch tán về phía sau.
        Giấu tin LSB thay thế bị kiểm định $\chi^2$ phát hiện ở $99.84\%$ ảnh stego và bị phá hủy bởi
        nén JPEG hay nhiễu nhẹ.
  \item \textbf{Về sản phẩm}: xây dựng ứng dụng desktop bốn tab chạy ngoại tuyến, gồm mã hóa và giấu
        tin, giải mã kèm phòng thí nghiệm tấn công, so sánh ba thuật toán, và bảng điều khiển kết quả
        8000 ảnh. Ứng dụng được đóng gói thành tệp \code{.exe} chạy độc lập trên Windows.
\end{itemize}

Về câu hỏi nên chọn thuật toán nào, kết quả của báo cáo cho thấy chất lượng ảnh stego \emph{không
phụ thuộc} thuật toán mã hóa, chỉ phụ thuộc số bit nhúng. Trong ba lựa chọn, \textbf{AES-128} là
phương án cân bằng nhất: nhanh nhờ AES-NI, là chuẩn được công nhận rộng rãi, và chống được tấn công
biết bản rõ. Phương án Chaotic + PRNG của bài báo chỉ nhanh hơn khi mặt nạ được tính sẵn, và không đủ
an toàn cho ứng dụng thực tế ở dạng hiện tại.

\section{Hạn chế}

\begin{itemize}
  \item Do bài báo không công bố mã nguồn, khóa và nhiều chi tiết cài đặt, bản tái hiện chỉ có thể
        bám sát ở mức các đại lượng có thể suy luận được. Một số chênh lệch nhỏ (SSIM, entropy ảnh
        stego) chưa xác định được nguyên nhân chính xác.
  \item Thời gian mã hóa được đo trên một máy tính cụ thể (Intel Core i3-1005G1) bằng Python, nên chỉ
        có ý nghĩa so sánh tương đối với MATLAB.
  \item Các cải tiến đề xuất ở \bang{tab:improve} mới dừng ở mức phân tích, chưa được cài đặt và đánh
        giá đầy đủ trong khuôn khổ đề tài.
  \item Phần đánh giá phân tích giấu tin mới dùng kiểm định $\chi^2$; chưa thử các phương pháp mạnh
        hơn như phân tích RS \cite{fridrich2001} hay các bộ phân loại học sâu.
\end{itemize}

\section{Hướng phát triển}

\begin{itemize}
  \item \textbf{Nâng cấp lớp mật mã}: thay CBC bằng mã hóa có xác thực AES-GCM \cite{sp80038d} với
        nonce ngẫu nhiên cho mỗi ảnh, đưa header vào dữ liệu liên kết để phát hiện mọi chỉnh sửa; dẫn
        xuất khóa từ mật khẩu bằng Argon2 \cite{rfc9106}.
  \item \textbf{Cải tiến lớp hỗn loạn}: thiết kế lại theo các yêu cầu của Alvarez và Li
        \cite{alvarez2006}, gồm khóa bí mật đủ dài, dòng khóa phụ thuộc nonce và có vòng khuếch tán
        trên toàn ảnh; sau đó đánh giá lại bằng phân tích vi sai đúng định nghĩa.
  \item \textbf{Cải tiến lớp giấu tin}: thử nghiệm LSB matching \cite{mielikainen2006}, hoán vị vị
        trí nhúng bằng khóa, giấu tin thích nghi theo độ phức tạp vùng ảnh \cite{holub2014}; đánh giá
        bằng phân tích RS và các bộ phát hiện dựa trên học máy.
  \item \textbf{Tăng độ bền vững}: kết hợp mã sửa lỗi và nhúng trong miền DCT/DWT để thông điệp sống
        sót qua nén JPEG.
  \item \textbf{Mở rộng ứng dụng}: hỗ trợ ảnh bí mật màu, video và âm thanh như hướng phát triển mà
        bài báo gốc đề xuất.
\end{itemize}
